\section{结论先行：利润是真，安全边际不足}

恒逸石化的投资问题不是“2026年利润是否反转”，而是这轮反转中有多少利润能够持续、归属于上市公司股东并转化为现金。公司预告2026H1归母净利润55--60亿元、扣非归母54.7--59.7亿元，经营改善具有较强可信度；但预告未经审计，也没有拆分文莱一期、国内PTA、聚酯及CPL--PA6的利润贡献。用半年利润机械年化会得到110--120亿元，这既忽略异常裂解价差均值回归，也忽略少数股东、利息、营运资金和潜在股份支付费用，因而不是可接受的估值分母。

本机构采用H1预告中值57.5亿元加H2正常化20.0亿元，得到House 2026E归母净利润77.5亿元。按已披露稀释股本38.80694189亿股计算，基准EPS约1.997元。模型不回加未量化的可转债有效利息；第七期员工持股计划按最多1.508138亿股、12.43元/股处理，并将最高18.746155亿元受让现金加入权益价值。由此得到基本面、市场和归一化Street三项锚15.6398/14.86/16.3130元，模型值15.6575元，发布目标取15.7元。相对2026-07-22 13:52 CST盘中价15.06元，上行约4.3\%，评级为\textbf{中性／观察}。

\begin{dashboardbox}[IC决策卡]
\noindent\begin{tabularx}{\textwidth}{L{2.8cm}X L{2.8cm}X}
\textbf{股票} & 恒逸石化（000703.SZ） & \textbf{当前价} & 15.06元（盘中） \\
\textbf{评级／研究状态} & 中性／观察 & \textbf{发布目标} & 15.7元 \\
\textbf{隐含空间} & 约+4.3\% & \textbf{模型值} & 15.6575元 \\
\textbf{核心公允区间} & 13.5--16.5元 & \textbf{尾部情景} & 10.5--22.5元 \\
\textbf{House 2026E归母} & 77.5亿元 & \textbf{已披露稀释EPS} & 约2.00元 \\
\textbf{下一结果桥} & H1中值57.5 + H2正常化20.0 & \textbf{核心风险} & 价差、现金流、杠杆、稀释 \\
\bottomrule
\end{tabularx}
\end{dashboardbox}
\sourcenote{HY-MKT-QUOTE（2026-07-22 13:52 CST）；HY-OFF-2026H1P（未经审计）；估值模型与审计，截止2026-07-22。}

这张决策卡的含义很直接：当前不是“看空反转”，而是“承认反转后仍要求价格纪律”。13.5元以下且基本面未证伪时，进入有条件研究升级区；13.5--16.3元维持观察并等待中报／三季报；16.3--16.5元降低研究优先级；16.5元以上若没有全年盈利与现金流同步上修，则触发重新评估。

\begin{keyinsight}[投资行动]
已知的H1利润已经被股价大幅交易，下一阶段回报取决于H2利润坡度和现金回款，而不是预告同比增速。\textbf{研究上，不因15元上方的静态“低PE”提高优先级；中报首先核验经营现金流、文莱分部和少数股东归属，三季报再核验Q3单季与9M累计是否接近12.0亿元和69.5亿元的House基准。}
\end{keyinsight}

\section{最终估值表与行动翻译}

本报告只对恒逸石化发布目标价。五家同行用于校验业务模式、周期位置和估值倍数，不进入本案核心估值池，因此不存在多股票排名。对单一标的的“排序”被转化为三项行动权重：盈利兑现与现金转化50\%、资本结构与稀释30\%、价格与事件位置20\%。盈利已经显著兑现，但现金、资本开支和稀释仍拖累综合分数，最终行动只能是观察而非积极增配。

\begin{exhibitbox}[最终估值、证据质量与失效条件]
\small
\noindent\begin{tabularx}{\textwidth}{L{3.35cm}R{2.45cm}X}
\toprule
\textbf{项目} & \textbf{数值} & \textbf{口径／证据} \\
\midrule
当前价／法定股本市值 & 15.06元／575.53亿元 & 38.21562147亿股法定股本；盘中行情 \\
House 2026E归母／EPS & 77.5亿元／1.997元 & H1中值57.5 + H2正常化20.0；约38.81亿股已披露稀释分母 \\
熊／基／牛价值 & 10.53／15.46／22.52元 & 65／77.5／95亿元归母，6.0／7.5／9.0倍周期PE，含员工计划匹配现金 \\
模型值／发布目标 & 15.6575／15.7元 & 基本面、市场、归一化Street三锚；相对当前价约+4.3\% \\
核心／尾部区间 & 13.5--16.5／10.5--22.5元 & 核心为House基准6.5--8.0倍；尾部为熊牛情景 \\
证据质量 & 官方+原始券商+盘中行情 & 财务、股本取官方；Street取完整原始PDF；概率与倍数为House判断 \\
评级／失效条件 & 中性／观察 & H2归母低于18亿元；现金流不修复；裂解显著回落；资本开支与融资压力上升 \\
\bottomrule
\end{tabularx}
\end{exhibitbox}
\sourcenote{HY-MKT-QUOTE；HY-OFF-2026H1P；HY-OFF-CB；HY-OFF-ESOP7；HY-BRK-GUOXIN；AStock House model，截止2026-07-22。}

这里的“中性”不是风险很低，而是预期收益与风险不对称。报告将“高风险”定义为：单一变量的合理不利变动即可令基准价值下移超过15\%，或需要公司外部融资才能维持既定扩张路径。按此定义，裂解价差、营运资金、杠杆与文莱二期资本开支均属于高风险项；目标价只给有限上行，并不意味着业务波动有限。

\section{从全链条到每股价值：利润池可以迁移，周期不会消失}

恒逸的资产覆盖原油输入、文莱炼油与芳烃、国内PTA/PIA、聚酯纤维与瓶片，以及CPL--PA6。链条完整的优势是利润池能够在炼油、芳烃和聚酯加工环节之间迁移；缺点是公司同时承担商品库存、少数股东、跨境物流、资本开支和高杠杆。下图因此不是“全产业链必然受益”的宣传图，而是一张利润必须经过归属和现金约束才能落到每股价值的地图。

\begin{exhibitbox}[全链条与估值池]
\centering
\begin{tikzpicture}[
  node distance=7mm and 6mm,
  block/.style={draw=rulegrey, rounded corners=2pt, fill=lightgrey, text width=2.35cm, align=center, minimum height=1.15cm, font=\small},
  core/.style={block, draw=navy, fill=blue!5},
  demandstyle/.style={block, draw=riskamber, fill=yellow!7},
  risk/.style={block, draw=riskred, fill=red!5},
  arrow/.style={-{Stealth[length=2mm]}, thick, draw=steelgrey}
]
\node[demandstyle] (oil) {原油／凝析油\\成本与库存锚};
\node[core, right=of oil] (brunei) {文莱一期\\炼油、PX、苯\\核心利润池};
\node[core, right=of brunei] (pta) {国内PTA／PIA\\参控股口径};
\node[core, right=of pta] (poly) {聚酯／瓶片\\体量最大};
\node[demandstyle, right=of poly] (demand) {织造、服装、包装\\需求锚};
\node[core, below=of pta] (nylon) {CPL／PA6\\事件期权};
\node[risk, below=of brunei] (constraints) {30\%少数股东\\利息、运费、税};
\node[risk, below=of poly] (cash) {库存、资本开支\\股份与现金};
\draw[arrow] (oil) -- (brunei);
\draw[arrow] (brunei) -- node[above,font=\scriptsize]{PX／苯与现金} (pta);
\draw[arrow] (pta) -- (poly);
\draw[arrow] (poly) -- (demand);
\draw[arrow] (nylon) -- (poly);
\draw[arrow] (constraints) -- (brunei);
\draw[arrow] (cash) -- (poly);
\end{tikzpicture}
\end{exhibitbox}
\sourcenote{HY-OFF-2025AR；HY-OFF-2026Q1；HY-OFF-P2；IND-CZCE；analysis/full\_chain\_taxonomy.md，截止2026-07-22。需求锚只表示终端强弱，不证明恒逸订单或收入。}

核心估值池只有恒逸石化。荣盛石化、恒力石化、东方盛虹、桐昆股份和新凤鸣属于卫星可比池；中国纺织服装、东南亚成品油和包装市场属于需求锚；付费Platts逐日价差、未具名客户、即时装置负荷等属于缺口池。覆盖缺口不会阻止对恒逸进行周期归一化估值，却阻止把2026H1峰值利润、文莱二期、CPL--PA6或瓶片认证给出无条件成长信用。

\section{紧凑供应链证据：什么能进入估值，什么不能}

读者需要区分“公司确有产品和资产”与“该资产在2026年赚了多少”。前者已经得到年报和项目公告支持；后者仍缺分部利润、实际价差、税率和客户映射。因此本报告用审计分部与情景估值，不用“设计产能乘行业价差”的伪精确方法。

\begin{exhibitbox}[公司级链条证据与缺口]
\small
\noindent\begin{tabularx}{\textwidth}{L{2.1cm}L{2.2cm}L{2.45cm}L{2.65cm}L{2.15cm}X}
\toprule
\textbf{链条角色} & \textbf{产品／工艺} & \textbf{客户／平台／应用证据} & \textbf{收入或规模证据} & \textbf{估值信用} & \textbf{缺失验证} \\
\midrule
文莱上游 & 800万吨炼化；油品565；PX／苯265万吨 & 产品面向东南亚和澳大利亚；客户未具名 & 2025主体收入421.40亿元、净利2.42亿元 & 周期情景可用 & 2026分部利润、实际价差、税率、客户与长协 \\
国内中游 & PTA／PIA & 对接聚酯；内部消化和外销量未拆 & PTA收入44.16亿元、毛利率0.37\%；PIA收入14.12亿元 & 审计分部可用 & 参控股权益、开工、转移价与单吨成本 \\
聚酯下游 & POY／FDY／DTY／短纤／切片 & 织造、服装、家纺、工业；客户匿名 & 涤纶收入597.44亿元、毛利率4.07\% & 合并周期信用 & 分品种量价、库存、开工、订单与ASP \\
包装材料 & 瓶片／RPET & 饮料、食品、包装为需求锚 & 产能530万吨，收入未单列 & 不给独立溢价 & 认证、客户、开工、收入与毛利 \\
尼龙链 & CPL／PA6 & 锦纶纺织与工程塑料应用 & 锦纶及副产品收入18.20亿元、毛利率6.83\% & 事件期权 & 良率、认证、分产品量价与单吨利润 \\
\bottomrule
\end{tabularx}
\end{exhibitbox}
\sourcenote{HY-OFF-2025AR（截至2025-12-31）；HY-OFF-2026H1P；客户链审计与供应链关系表，研究截止2026-07-22。}

这张表对投资行动的约束是：文莱一期和国内合并业务可以进入基准，但未披露的客户、订单、瓶片认证、CPL--PA6良率及二期融资只能进入监控清单。产业链映射提高了理解力，却不能绕过财务披露门槛。

\section{价值链经济性：利润最强处与现金最弱处并存}

2026H1最强利润池位于文莱成品油和芳烃；国内PTA从2025年低加工费阶段性修复，聚酯仍受织机开工和长丝库存压制；CPL--PA6处于爬坡阶段。公司一季度利润率上升的同时，存货和经营现金流恶化，说明价值池的会计兑现与现金兑现不同步。

\begin{exhibitbox}[价值链经济性与估值信用]
\small
\noindent\begin{tabularx}{\textwidth}{L{1.85cm}L{2.25cm}L{2.25cm}L{2.0cm}L{2.1cm}X}
\toprule
\textbf{环节} & \textbf{价值／价格代理} & \textbf{利润池} & \textbf{供需／负荷} & \textbf{订单／认证} & \textbf{估值处理} \\
\midrule
原油输入 & Brent、采购贴水、运费、库存损益 & 原料约占炼油成本96.2\% & 原油供给与产品紧张并存 & 不适用 & 仅作成本情景，不把油价涨跌直接写成利润 \\
文莱油品／芳烃 & 柴油／航煤裂解、PX--石脑油价差 & 2026H1最强 & 公司称满产满销 & 客户与长协未披露 & 纳入熊基牛情景，不把高景气永久化 \\
PTA & $PTA-0.655\times PX$ & 2025平均加工差255.6元／吨；2026年6月代理665元／吨 & 检修去库后阶段修复 & 商品业务不以backlog估值 & 只用分部及情景；禁止2150万吨乘加工费 \\
聚酯 & 加工费、开工、库存 & 2025毛利率4.07\% & 2026年6月行业库存仍高 & 客户和分品种认证未披露 & 合并周期估值；不给分品种溢价 \\
CPL／PA6 & 单吨价差尚未取得 & 爬坡期权 & 2025产8.89、销4.67、存4.22万吨 & 良率／认证／客户未知 & 基准零增量信用，等半年报验证 \\
资本与现金 & OCF、库存、利息、资本开支 & 决定利润能否沉淀 & Q1 OCF为-25.73亿元 & 不适用 & 压低倍数；二期基准零信用 \\
\bottomrule
\end{tabularx}
\end{exhibitbox}
\sourcenote{HY-OFF-2025AR；HY-OFF-2026Q1；IND-PTA-202606；IND-CZCE；analysis/value\_chain\_economics.md，截止2026-07-22。行业代理不等于恒逸单吨利润。}

投资含义是：基本面上行与现金风险可以同时存在。只要文莱裂解和负荷强，利润仍有支撑；但若H1经营现金流未显著修复、库存和短债继续上升，就应先下调估值倍数，而不是继续上修峰值利润。

\section{盈利、预期与估值三座桥}

下一季度不是再预测一次同比增速，而是完成从预告到审计、从总利润到归母、从利润到现金的三层核验。House 2026E基准77.5亿元略低于国信79.1亿元，但差异不大；真正的差异在于House明确让H2利润降到20亿元，并在分母中加入可转债潜在股份和员工计划上限。

\begin{exhibitbox}[2026E预期—估值桥]
\noindent\begin{tabularx}{\textwidth}{L{3.7cm}R{2.1cm}X}
\toprule
\textbf{项目} & \textbf{数值} & \textbf{解释} \\
\midrule
2026E营业收入 & 1,417.4亿元 & 采用可审计国信原始PDF作规模检查，不是公司指引 \\
收入增速 & 约24.8\% & 相对2025A 1,135.27亿元的模型检查值 \\
House 2026E归母／EPS & 77.5亿元／1.997元 & H1中值57.5 + H2正常化20.0；EPS按38.80694189亿股 \\
基准预期倍数 & 7.5倍PE & 周期调整基准，不采用结构性成长PEG \\
基准情景权益价值 & 15.4611元 & $[77.5\times7.5+18.746155]\div38.80694189$ \\
基本面概率锚 & 15.6398元 & 熊／基／牛25\%／55\%／20\%加权 \\
最终模型／发布目标 & 15.6575／15.7元 & 再加入市场锚15\%与Street锚20\% \\
\bottomrule
\end{tabularx}
\end{exhibitbox}
\sourcenote{HY-OFF-2025AR；HY-OFF-2026H1P；HY-BRK-GUOXIN；analysis/earnings\_forecast.md与analysis/valuation\_audit.md，截止2026-07-22。}

2026E收入与归母预测并非来自同一权重来源：收入采用国信原始报告作总量检查，House归母则是独立的H1--H2正常化判断。这样可以避免把79.1亿元同时放进基本面锚和Street锚，造成隐性双重计权。

\begin{exhibitbox}[市场隐含盈利：股本口径决定“便宜”的程度]
\noindent\begin{tabularx}{\textwidth}{L{4.1cm}R{2.0cm}R{2.0cm}X}
\toprule
\textbf{股本状态} & \textbf{隐含归母} & \textbf{占House基准} & \textbf{用途} \\
\midrule
当前经济股本34.3972亿股 & 69.07亿元 & 89.1\% & 未转股、未回流库存股的交易视角 \\
加入转债潜在股份37.2988亿股 & 74.90亿元 & 96.6\% & 不回加有效利息的保守稀释敏感性 \\
已披露稀释38.8069亿股并扣员工现金 & 75.42亿元 & 97.3\% & 主估值最接近的每股敏感性 \\
\bottomrule
\end{tabularx}
\end{exhibitbox}
\sourcenote{HY-MKT-QUOTE；HY-OFF-CB；HY-OFF-DIV；HY-OFF-ESOP7；analysis/valuation\_audit.md，截止2026-07-22。后两项为稀释敏感性，不标作会计if-converted EPS。}

这组反推显示，所谓“7倍PE很便宜”在完全稀释后不再形成很厚的安全边际。当前价格已经计入House基准的约97\%，投资者真正押注的是基准利润上修、倍数扩张或牛情景概率上升。

\section{三锚目标、Street比较与最强反方论证}

本报告让基本面锚占65\%，市场锚占15\%，Street锚占20\%。基本面权重最高，因为官方财务和股本可复核；市场锚限制模型脱离交易现实；Street只有一份原始PDF可获正权重，其余七份媒体转载摘要权重为零。这不是“一致预期均值”，而是一份可审计外部锚与公开摘要温度的分层使用。

\begin{exhibitbox}[市场情绪锚与最终权重]
\noindent\begin{tabularx}{\textwidth}{L{3.2cm}R{1.8cm}R{1.6cm}X}
\toprule
\textbf{锚} & \textbf{每股价值} & \textbf{最终权重} & \textbf{证据与作用} \\
\midrule
House基本面概率锚 & 15.6398元 & 65\% & 独立77.5亿元归母、三情景PE、员工计划现金 \\
市场隐含锚 & 14.86元 & 15\% & 五日VWAP 14.6555元与盘中价15.06元中点 \\
归一化Street锚 & 16.3130元 & 20\% & 国信目标中值17.05元按股本与员工现金归一 \\
最终模型值 & 15.6575元 & 100\% & $15.6398\times65\%+14.86\times15\%+16.3130\times20\%$ \\
发布目标 & 15.7元 & -- & 对详细模型小数作可读化处理 \\
\bottomrule
\end{tabularx}
\end{exhibitbox}
\sourcenote{HY-MKT-KLINE；HY-MKT-QUOTE；HY-BRK-GUOXIN；analysis/valuation\_audit.md，截止2026-07-22。}

\begin{exhibitbox}[公开Street比较与证据质量]
\small
\noindent\begin{tabularx}{\textwidth}{L{2.0cm}L{1.7cm}L{1.5cm}L{2.0cm}L{2.4cm}X}
\toprule
\textbf{券商／来源} & \textbf{日期} & \textbf{评级} & \textbf{目标／预测} & \textbf{方法} & \textbf{证据质量与权重} \\
\midrule
国信证券 & 2026-07-04 & 优于大市 & 16.5--17.6元；2026E归母79.1亿元 & FCFE + 可比PE & 原始PDF；正权重20\% \\
光大／中金／国海／中信建投等7份 & 2026-01-06至06-28 & 买入／跑赢 & 摘要可见字段不完整 & 未取得完整方法 & 媒体转载摘要；估值权重0 \\
AStock House & 2026-07-22 & 中性／观察 & 目标15.7元；2026E归母77.5亿元 & 周期PE + 市场 + 归一化Street & 官方数据与模型复算；报告主判断 \\
\bottomrule
\end{tabularx}
\end{exhibitbox}
\sourcenote{HY-BRK-GUOXIN；data/broker\_street\_consensus\_20260722.md；sources/broker-reports/2026-07-22/index.md，检索截止2026-07-22。}

最强反方观点是：House把H2正常化得过快。如果东南亚成品油短缺、炼厂供给收缩和PX／苯价差在三季度继续维持高位，文莱一期保持满负荷，国内PTA加工费与聚酯库存又同步改善，那么H2归母可以接近牛情景37.5亿元，95亿元全年归母和9倍PE将支持约22.5元价值。相反，House错误的证据不是股价上涨，而是半年报和三季报同时证明高裂解价差、现金覆盖利息与资本开支、国内库存下降且二期融资受控。

\begin{riskbox}[什么会证明本报告错了]
若2026H2归母达到35亿元以上，H1经营现金流显著转正并覆盖利息和资本开支，文莱分部披露能够解释少数股东与税惠，国内PTA／聚酯库存和加工费同步改善，则当前“中性／观察”过于保守，应把牛情景提高为主情景。\textbf{研究上，只有这些证据组合出现，才把16.5元以上从重新评估区改为基本面上修区。}
\end{riskbox}

反过来，若H2归母低于18亿元、经营现金流仍大额为负，或库存、短债、在建工程同时上升，则不仅利润预测要下调，基准PE也应从7.5倍向6倍收敛。尾部10.5元并非“极端灾难价”，而是高景气回落与资本约束同时发生时可以复算的熊情景。

\section{股本与员工计划：本报告最容易被误读的两项边界}

第一，剩余2.31亿股库存股用途尚未确认。它们已经从当前经济股本中扣除，但不能假设永久注销，也不能在没有公告时预设用于下一期员工计划、出售或其他资本运作。若未来回流，EPS分母和现金流入必须同时更新。第二，第七期员工计划披露将涉及股份支付会计处理，但公司没有给出2026年精确费用。主模型使用报告口径利润，不自行回加股份支付，也不把未知费用填成零风险。

\begin{exhibitbox}[关键资本结构边界]
\noindent\begin{tabularx}{\textwidth}{L{3.8cm}R{2.2cm}X}
\toprule
\textbf{项目} & \textbf{数量／金额} & \textbf{主模型处理} \\
\midrule
法定总股本 & 38.21562147亿股 & 官方时点股本，不直接作为EPS分母 \\
回购库存股 & 3.81840072亿股 & 从经济股本扣除 \\
剩余转债潜在新增 & 2.90158314亿股 & 加入已披露稀释分母；不回加未量化有效利息 \\
第七期员工计划上限 & 1.508138亿股 & 加入稀释分母 \\
员工计划受让价／最高现金 & 12.43元／18.746155亿元 & 现金加入权益价值，最终认购不确定 \\
已披露稀释股本 & 38.80694189亿股 & 主估值分母 \\
其余库存股 & 2.31026272亿股 & 用途未确认，列为边界风险，不假设永久注销 \\
股份支付费用 & 公司未量化 & 不回加、不伪造；中报后更新 \\
\bottomrule
\end{tabularx}
\end{exhibitbox}
\sourcenote{HY-OFF-CB（2026-06-30股本与转债）；HY-OFF-DIV（回购专户分红登记口径）；HY-OFF-ESOP7（2026-06-13计划上限、价格与会计处理）；analysis/valuation\_audit.md。}

因此，本报告的4.3\%上行不是对员工计划“利好”或“利空”的单边判断，而是在最大股份回流与最大现金流入同时发生的披露上限情景下得到。最终认购量、过户时点和股份支付费用一旦披露，必须同时更新股本、现金和利润三项，不能只挑对目标价有利的一边。
